\chapter{Charges and population analysis}
\label{ch:charges}

Once a density has been predicted, two questions follow: how much charge is
there, and where does it sit. This chapter covers the conservation check that
must pass before either answer means anything, and the three partitionings that
turn a grid into per-atom numbers.

Nothing here is learned. Every quantity is a deterministic functional of a
density that has already been predicted, so \textbf{the error in a partial
charge is inherited from the density and is never smaller than it}.

\section{Charge conservation}
\label{sec:charge-conservation}

A Kohn--Sham density integrates to the valence count the pseudopotentials fix.
That number is an \emph{input} to a DFT calculation rather than an output of
one, so it is exactly known --- and it is the only property of a predicted
density that can be checked without a reference calculation.

\begin{lstlisting}[language=Python]
from poraque.analysis import verify_total_charge

check = verify_total_charge(rho.data, rho.grid.cell, expected_electrons=297.0)
print(check)
# charge check OK: 297.000001 electrons against an expected 297.000000
#                  (+1.918e-09 relative, tolerance 0.001)
\end{lstlisting}

The integral is the voxel sum times the voxel volume,

\begin{equation*}
\int\rho\,d^3r = \sum_{ijk}\rho_{ijk}\,\mathrm{d}V,
\qquad
\mathrm{d}V = \frac{\Omega}{N_x N_y N_z},
\end{equation*}

which is \emph{exact} for a band-limited field on a uniform periodic mesh: the
rectangle and trapezoid rules coincide under periodic boundary conditions, so
there is no quadrature error to argue about, only the field itself.

The default tolerance of $10^{-3}$ is not arbitrary. The electrostatic terms
are of order $10^{4}$ eV, so a relative drift of $10^{-3}$ already moves a
total energy by roughly $10$ eV --- more than any energy difference worth
computing.

\begin{pwarn}
The shipped \opt{ext2chg} operator predicts densities whose electron count
drifts by up to $1.7\,\%$ --- nearly five electrons out of $297$. Left
uncorrected this moves totals by tens of eV, differently for each structure, so
it does not cancel in a difference.
\end{pwarn}

\subsection{Automatic normalisation}

The calculator rescales the density to the exact count by default:

\begin{lstlisting}[language=Python]
atoms.calc = Poraque("models/poraque_models.pfno", potcar="POTCAR",
                     normalize_density=True)   # the default

atoms.calc.verify_charge()          # passes by construction
atoms.calc.raw_electron_drift       # what the operator actually did
\end{lstlisting}

With normalisation on, the check passes by construction, so the informative
number afterwards is \opt{raw\_electron\_drift} --- the drift of the raw
prediction, recorded before the repair. Set \opt{normalize\_density=False} to
measure the operator rather than the repair.

Rescaling is a repair and not a cure: a density whose integral is wrong by
$1.7\,\%$ has the wrong shape too, and a single global factor does not fix a
shape.

\section{Partial charges}

All three schemes share one form,

\begin{equation*}
q_A = Z^{\rm val}_A - \int w_A(\rr)\,\rho(\rr)\,d^3r,
\qquad
\sum_A w_A(\rr) = 1,
\end{equation*}

and differ only in the weight $w_A$. Because the partitions are exhaustive,
\textbf{all three conserve charge exactly}: the populations sum to
$\int\rho$ whatever the weights are.

\begin{center}
\begin{tabular}{@{}llll@{}}
\toprule
Method & $w_A(\rr)$ & Cost & Character \\
\midrule
\opt{voronoi} & 1 for the nearest atom & fast & purely geometric \\
\opt{hirshfeld} & $\rho^{\rm at}_A / \sum_B \rho^{\rm at}_B$ & fast & needs a promolecule \\
\opt{bader} & 1 inside $A$'s zero-flux basin & moderate & intrinsic to $\rho$ \\
\bottomrule
\end{tabular}
\end{center}

\subsection{From the ASE calculator}

\begin{lstlisting}[language=Python]
atoms.calc = Poraque("models/poraque_models.pfno", potcar="POTCAR",
                     references="data/vasp/ref")

atoms.get_charges()                       # standard ASE, Bader by default
atoms.calc.get_charges(method="hirshfeld")
atoms.calc.get_charges(method="voronoi")

print(atoms.calc.charge_analysis)         # the full decomposition
\end{lstlisting}

\opt{get\_charges} returns net charges in units of $+e$, positive for
electron-deficient, and is compatible with the standard
\opt{atoms.get\_charges()} interface. The full analysis --- populations, the
valence subtracted, which promolecule or Bader backend was used --- is left on
\opt{charge\_analysis}.

\subsection{Voronoi}

Each voxel goes wholly to its nearest atom under the minimum-image convention.
In a non-orthogonal cell the surrounding images are searched rather than
trusting a wrap of the fractional coordinates --- in a skewed lattice the
nearest image can be a diagonal neighbour.

Voxels exactly equidistant from several atoms are \textbf{shared equally}
rather than awarded to the lowest atom index. That is rare on a real density
and common on a symmetric cell sampled by a commensurate grid, which is exactly
where an index-order tie-break would bias a result that ought to come out
symmetric.

Purely geometric: the density enters only through the integral, never through
where the boundary is drawn. For atoms of unequal size that is a poor chemical
partition, which is why it is a baseline rather than a default.

\subsection{Hirshfeld}

Weights by the free atoms, so the partition is smooth and chemically sensible
--- but it is \emph{defined by its reference}, and a poor promolecule gives
poor charges with no other symptom.

\begin{lstlisting}[language=Python]
partial_charges(rho, method="hirshfeld", valence={"Au": 11.0},
                references="data/vasp/ref")
\end{lstlisting}

Each \file{<references>/<Element>/CHGCAR} is spherically averaged into
$\rho^{\rm at}(r)$. Where no reference exists the fallback is an exponential
free atom, $\rho(r) = \frac{N}{8\pi a^3}e^{-r/a}$, normalised to the valence
count with the decay length set from the covalent radius through
$\langle r\rangle = 3a$. Which was used for each element is reported in
\opt{details["promolecule"]} --- check it, because the fallback is crude.

\begin{pnote}
Hirshfeld charges are small in magnitude compared with other schemes. For an
elemental solid they come out near zero by construction, since the promolecule
is built from identical free atoms; that makes a useful sanity check on the
reference handling.
\end{pnote}

\subsection{Bader (QTAIM)}

The only scheme whose definition is intrinsic to the density: basins are
bounded by surfaces of zero flux in $\nabla\rho$. Two backends are provided.

\begin{lstlisting}[language=Python]
partial_charges(rho, method="bader", backend="native")    # pure NumPy
partial_charges(rho, method="bader", backend="external")  # Henkelman `bader`
partial_charges(rho, method="bader", backend="auto")      # external if present
\end{lstlisting}

\textbf{Native} is a vectorised on-grid steepest ascent: every voxel points at
whichever of its 26 neighbours gives the largest density increase \emph{per
unit distance}, and those pointers are followed to a local maximum by repeated
pointer doubling. The distance weighting matters --- without it the diagonal
neighbours, $\sqrt3$ further away, win too often and the basins acquire a
staircase bias along the cell diagonals.

\textbf{External} writes a temporary \file{CHGCAR}, runs the Henkelman group's
\opt{bader} program and parses \file{ACF.dat}. \opt{backend="external"} raises
if the program is absent rather than falling back silently: the two backends
are not identical, and a result reported as ``Bader (Henkelman)'' should not
quietly have been something else.

A density may have more maxima than atoms --- spurious ones from grid noise, or
genuine non-nuclear attractors in a metal. Each maximum is assigned to its
nearest atom, so the mapping is many-to-one, and \opt{details["maxima"]}
reports how many were found.

\section{What these charges are not}

\begin{pwarn}
All three partition the \emph{pseudo} valence density. The PAW core is absent,
so these are not all-electron populations: a Bader volume here is not the
all-electron Bader volume, and the charges are systematically compressed toward
zero. This matters most for Bader, whose basin boundaries are set by the
density's topology near the nuclei --- exactly where the pseudisation is
largest.

For a proper Bader analysis VASP's own guidance is to sum \file{AECCAR0} and
\file{AECCAR2}; Poraqu\^e does not predict those.

Treat the numbers as comparative across a series, not as absolute.
\end{pwarn}

\section{In the reports}

Passing a \pyclass{PartialCharges} to the report builder typesets the per-atom
table, the method and the caveat above into the PDF:

\begin{lstlisting}[language=Python]
from poraque.vis import ModelReport

ModelReport("reports").build(task="ext2chg", per_material=metrics,
                             charges=analysis)
\end{lstlisting}
